\chapter{Clustering Methods}
\label{chap:methods}

The methodology adopted in this research is divided into three main stages: the mathematical definition of the clustering algorithms selected for comparison, the generation of synthetic data using the CART method, and the simulation pipeline designed to evaluate the structural degradation of the data.

\section{Clustering Algorithms}
\label{sec:clustering-algorithms}

Clustering is the task of grouping a set of objects in such a way that objects in the same group (cluster) are more similar to each other than to those in other groups. This research focuses on two of the most established paradigms in unsupervised learning: centroid-based clustering (K-Means) and connectivity-based clustering (Hierarchical).

\subsection{K-Means Clustering}
K-Means is a partition-based algorithm that aims to divide $N$ observations into $k$ non-overlapping clusters. It operates by minimizing the within-cluster sum of squares (WCSS), also known as inertia \cite{macqueen1967}.

\textbf{Formal Mathematical Definition}

We define the K-Means procedure by the tuple
\begin{equation}
  \mathcal{M}_{KM} = \langle \mathcal{X}, k, \mathcal{P}, \boldsymbol{\mu}, \mathcal{A}, \mathcal{U}, J \rangle
\end{equation}
where each component is defined as follows:

\begin{itemize}
  \item \textbf{Data Space ($\mathcal{X}$):} The observed sample is
  \begin{equation}
    \mathcal{X} = \{\mathbf{x}_1, \mathbf{x}_2, \dots, \mathbf{x}_N\}, \qquad \mathbf{x}_i \in \mathbb{R}^d
  \end{equation}
  with $N$ observations and feature dimension $d$.

  \item \textbf{Number of Clusters ($k$):} A fixed integer satisfying
  \begin{equation}
    k \in \mathbb{N}, \qquad 1 < k < N
  \end{equation}
  that specifies the target cardinality of the partition.

  \item \textbf{Partition Space ($\mathcal{P}$):} A feasible clustering is a partition
  \begin{equation}
    \mathcal{P} = \{S_1, S_2, \dots, S_k\}
  \end{equation}
  such that
  \begin{equation}
    S_a \cap S_b = \varnothing \quad (a \neq b), \qquad \bigcup_{j=1}^{k} S_j = \mathcal{X}, \qquad S_j \neq \varnothing
  \end{equation}

  \item \textbf{Centroid Set ($\boldsymbol{\mu}$):} Each cluster $S_j$ is represented by a centroid
  \begin{equation}
    \boldsymbol{\mu} = \{\boldsymbol{\mu}_1, \dots, \boldsymbol{\mu}_k\}, \qquad \boldsymbol{\mu}_j \in \mathbb{R}^d
  \end{equation}
  with
  \begin{equation}
    \boldsymbol{\mu}_j = \frac{1}{|S_j|}\sum_{\mathbf{x}_i \in S_j} \mathbf{x}_i
  \end{equation}

  \item \textbf{Assignment Operator ($\mathcal{A}$):} For fixed centroids $\boldsymbol{\mu}^{(t)}$ at iteration $t$, the assignment map is
  \begin{equation}
    \mathcal{A}_{\boldsymbol{\mu}^{(t)}} : \mathcal{X} \to \{1,\dots,k\}, \qquad
    \mathcal{A}_{\boldsymbol{\mu}^{(t)}}(\mathbf{x}_i) = \arg\min_{1 \le j \le k} \|\mathbf{x}_i - \boldsymbol{\mu}^{(t)}_j\|_2^2
  \end{equation}
  and therefore induces the partition
  \begin{equation}
    S_j^{(t)} = \left\{\mathbf{x}_i \in \mathcal{X} : \mathcal{A}_{\boldsymbol{\mu}^{(t)}}(\mathbf{x}_i) = j \right\}
  \end{equation}

  \item \textbf{Update Operator ($\mathcal{U}$):} For a fixed partition $\mathcal{P}^{(t)} = \{S_1^{(t)}, \dots, S_k^{(t)}\}$, the centroid update is
  \begin{equation}
    \mathcal{U}_{\mathcal{P}^{(t)}}(S_j^{(t)}) = \boldsymbol{\mu}^{(t+1)}_j
    = \frac{1}{|S_j^{(t)}|}\sum_{\mathbf{x}_i \in S_j^{(t)}} \mathbf{x}_i, \qquad j = 1,\dots,k
  \end{equation}

  \item \textbf{Objective Function ($J$):} The algorithm minimizes the within-cluster sum of squares
\end{itemize}
\begin{equation}
  J(\mathcal{P}, \boldsymbol{\mu}) = \sum_{j=1}^{k}\sum_{\mathbf{x}_i \in S_j}\|\mathbf{x}_i - \boldsymbol{\mu}_j\|_2^2
\end{equation}
which is equivalently the inertia of the partition.

The iterative state of the algorithm can be written as
\begin{equation}
  \left(\mathcal{P}^{(t)}, \boldsymbol{\mu}^{(t)}\right) \mapsto
  \left(\mathcal{P}^{(t+1)}, \boldsymbol{\mu}^{(t+1)}\right)
\end{equation}
with the alternating minimization steps
\begin{equation}
  \mathcal{P}^{(t+1)} = \arg\min_{\mathcal{P}} J(\mathcal{P}, \boldsymbol{\mu}^{(t)}),
  \qquad
  \boldsymbol{\mu}^{(t+1)} = \arg\min_{\boldsymbol{\mu}} J(\mathcal{P}^{(t+1)}, \boldsymbol{\mu})
\end{equation}

The convergence criterion can be written as either
\begin{equation}
  \mathcal{P}^{(t+1)} = \mathcal{P}^{(t)}
  \qquad \text{or} \qquad
  \|\boldsymbol{\mu}^{(t+1)} - \boldsymbol{\mu}^{(t)}\|_F \le \varepsilon
\end{equation}
for some tolerance $\varepsilon > 0$, where $\|\cdot\|_F$ denotes the Frobenius norm.

\textbf{Time and Space Complexity}

Let $I$ denote the number of iterations until convergence. The main computational blocks are:
\begin{equation}
  T_{\text{init}} = O(kd)
\end{equation}
for selecting or storing the initial centroids,
\begin{equation}
  T_{\text{assign}}^{(t)} = O(Nkd)
\end{equation}
for computing the distances from all $N$ points to all $k$ centroids at iteration $t$, and
\begin{equation}
  T_{\text{update}}^{(t)} = O(Nd)
\end{equation}
for recomputing all centroids once the assignments are known.

Hence the total running time is
\begin{equation}
  T_{KM}(N,k,d,I) = T_{\text{init}} + \sum_{t=1}^{I}\left(T_{\text{assign}}^{(t)} + T_{\text{update}}^{(t)}\right)
  = O(kd) + O\!\left(I(Nkd + Nd)\right)
\end{equation}
which is dominated by
\begin{equation}
  T_{KM}(N,k,d,I) = O(NkId)
\end{equation}

The corresponding memory requirement is
\begin{equation}
  S_{KM}(N,k,d) = O(Nd) + O(kd) + O(N)
\end{equation}
for the data matrix, the centroid matrix, and the assignment vector, which simplifies to
\begin{equation}
  S_{KM}(N,k,d) = O(Nd + kd)
\end{equation}

\textbf{Practical Considerations}

K-Means is highly efficient in practice, but it has significant limitations relevant to this study:
\begin{itemize}
  \item It assumes clusters are spherical and of similar size.
  \item It is sensitive to outliers and the initial placement of centroids.
  \item The number of clusters $k$ must be specified \textit{a priori}.
\end{itemize}

\subsection{Hierarchical Clustering (Agglomerative)}
Unlike K-Means, Hierarchical Clustering does not require a pre-specified number of clusters. Instead, it builds a hierarchy of clusters, often represented as a dendrogram. This research utilizes \textbf{Agglomerative Clustering} (bottom-up), which treats each observation as its own cluster and successively merges pairs of clusters \cite{ward1963}.

\textbf{Formal Mathematical Definition (Ward's Linkage)}

Agglomerative Hierarchical Clustering with Ward's linkage is a deterministic, greedy, bottom-up agglomerative framework. It can be formally defined by the tuple
\begin{equation}
  \mathcal{M}_{HC} = \langle \mathcal{X}, \mathcal{P}^{(0)}, \mathcal{T}, \Delta, \mathcal{L}, \mathcal{M} \rangle
\end{equation}
where:

\begin{itemize}
  \item \textbf{Data Space ($\mathcal{X}$):} An input set of $N$ observations
  \begin{equation}
    \mathcal{X} = \{\mathbf{x}_1, \dots, \mathbf{x}_N\} \subset \mathbb{R}^d
  \end{equation}

  \item \textbf{Initial Partition ($\mathcal{P}^{(0)}$):} The algorithm starts from the singleton partition
  \begin{equation}
    \mathcal{P}^{(0)} = \left\{ \{\mathbf{x}_1\}, \{\mathbf{x}_2\}, \dots, \{\mathbf{x}_N\} \right\}
  \end{equation}

  \item \textbf{Hierarchy ($\mathcal{T}$):} A dendrogram represented as a sequence of nested partitions
  \begin{equation}
    \mathcal{T} = \{\mathcal{P}^{(0)}, \mathcal{P}^{(1)}, \dots, \mathcal{P}^{(N-1)}\}
  \end{equation}
  satisfying
  \begin{equation}
    |\mathcal{P}^{(t)}| = N - t, \qquad t = 0,1,\dots,N-1
  \end{equation}
  and each $\mathcal{P}^{(t+1)}$ is obtained by merging exactly two blocks from $\mathcal{P}^{(t)}$.

  \item \textbf{Error Sum of Squares ($\Delta$):} The spatial dispersion within any candidate cluster $C \subset \mathcal{X}$ is quantified by the Error Sum of Squares (ESS), defining the function $\Delta(C)$:
  \begin{equation}
    \Delta(C) = \text{ESS}(C) = \sum_{\mathbf{x} \in C} ||\mathbf{x} - \boldsymbol{\mu}_C||_2^2 \quad \text{where} \quad \boldsymbol{\mu}_C = \frac{1}{|C|}\sum_{\mathbf{x} \in C} \mathbf{x}
  \end{equation}
  \item \textbf{Linkage Objective Function ($\mathcal{L}$):} The merging cost between two disjoint clusters $A$ and $B$ is defined by Ward's minimal variance criterion. The algorithm greedily selects the pair $(A^*, B^*)$ that minimizes the incremental increase in the total intra-cluster variance:
  \begin{equation}
    \mathcal{L}(A, B) = \Delta(A \cup B) - \left[\Delta(A) + \Delta(B)\right] = \frac{|A||B|}{|A| + |B|} ||\boldsymbol{\mu}_A - \boldsymbol{\mu}_B||_2^2
  \end{equation}
  At each agglomerative step $t$, the transition $\mathcal{P}^{(t-1)} \rightarrow \mathcal{P}^{(t)}$ operates via:
  \begin{equation}
      (A_t^*, B_t^*) = \arg\min_{A, B \in \mathcal{P}^{(t-1)}, A \neq B} \mathcal{L}(A, B)
  \end{equation}
  and the merged partition becomes
  \begin{equation}
    \mathcal{P}^{(t)} =
    \left(\mathcal{P}^{(t-1)} \setminus \{A_t^*, B_t^*\}\right) \cup \{A_t^* \cup B_t^*\}
  \end{equation}

  \item \textbf{Merge Operator ($\mathcal{M}$):} The agglomeration rule may be written as
  \begin{equation}
    \mathcal{M}(\mathcal{P}^{(t-1)}) = \mathcal{P}^{(t)}
  \end{equation}
  so that the full algorithm is the recursive application of $\mathcal{M}$ for $t=1,\dots,N-1$.
\end{itemize}

\textbf{Equivalent Lance--Williams Update for Ward's Linkage}

If $A$ and $B$ are merged into $A \cup B$, then for any other active cluster $C$, the updated Ward distance can be written in Lance--Williams form as
\begin{multline}
  d_{W}(A \cup B, C)
  =
  \frac{|A| + |C|}{|A| + |B| + |C|}d_W(A,C)
  +
  \frac{|B| + |C|}{|A| + |B| + |C|}d_W(B,C) \\
  -
  \frac{|C|}{|A| + |B| + |C|}d_W(A,B)
\end{multline}
which shows explicitly how the pairwise merge costs are updated after each agglomerative step.

\textbf{Time and Space Complexity}

Let $m_t = N - t + 1$ denote the number of active clusters at step $t$. In a naive implementation, the dominant tasks are:
\begin{equation}
  T_{\text{dist}} = O(N^2 d)
\end{equation}
to construct the initial pairwise dissimilarity structure,
\begin{equation}
  T_{\text{search}}^{(t)} = O(m_t^2)
\end{equation}
to search for the minimum merge cost among all active cluster pairs, and
\begin{equation}
  T_{\text{update}}^{(t)} = O(m_t)
\end{equation}
to update the linkage values after a merge.

Therefore, the total cost can be expressed as
\begin{equation}
  T_{HC}(N,d)
  =
  O(N^2 d)
  +
  \sum_{t=1}^{N-1}\left(O(m_t^2) + O(m_t)\right)
\end{equation}
and since $\sum_{t=1}^{N-1} m_t^2 = O(N^3)$, this yields
\begin{equation}
  T_{HC}(N,d) = O(N^2 d + N^3) = O(N^3)
\end{equation}
in the standard naive formulation.

With optimized data structures, the practical running time can be reduced to approximately
\begin{equation}
  T_{HC}^{\text{opt}}(N,d) = O(N^2 \log N) \quad \text{or} \quad O(N^2)
\end{equation}
depending on the specific implementation.

The storage requirement is dominated by the pairwise dissimilarity representation:
\begin{equation}
  S_{HC}(N,d) = O(Nd) + O(N^2)
\end{equation}
which simplifies asymptotically to
\begin{equation}
  S_{HC}(N,d) = O(N^2)
\end{equation}

\textbf{Practical Considerations}

Hierarchical clustering is computationally expensive, making it less suitable for massive datasets. However, it offers distinct advantages:
\begin{itemize}
  \item It captures hierarchical structures and does not assume strictly spherical clusters.
  \item It is deterministic (produces the same result every time), unlike K-Means which depends on random initialization.
  \item It allows for the determination of $k$ \textit{post-hoc} by cutting the dendrogram at a specific height.
\end{itemize}

\subsection{Comparative and Complexity Summary}

When selecting a clustering algorithm for synthetic data evaluation, understanding the computational constraints is as critical as understanding the geometric assumptions.

\textbf{Time Complexity:} The asymptotic contrast between the two methods can be summarized as
\begin{equation}
  T_{KM}(N,k,d,I) = O(NkId)
  \qquad \text{versus} \qquad
  T_{HC}(N,d) = O(N^3)
\end{equation}
for the standard forms of the algorithms. Since $k$, $I$, and $d$ are usually much smaller than $N$ in our simulation settings, K-Means is often treated as effectively linear in $N$, whereas Hierarchical Clustering remains superlinear because it must maintain and update inter-cluster dissimilarities throughout the agglomerative process.

\textbf{Computational (Space) Complexity:} The corresponding memory requirements are
\begin{equation}
  S_{KM}(N,k,d) = O(Nd + kd)
  \qquad \text{and} \qquad
  S_{HC}(N,d) = O(N^2)
\end{equation}
so K-Means scales linearly with the sample representation, while Hierarchical Clustering is dominated by the quadratic cost of storing and updating pairwise distances. This creates the main computational tradeoff between the two methods.

Table~\ref{tab:kmeans_vs_hc} summarizes the key computational, structural, and operational differences between the two algorithms evaluated in this study.

\begin{table}[H]
  \centering
  \caption[Comparison of K-Means and Hierarchical Clustering]{Comparison of K-Means and Hierarchical Clustering}
  \label{tab:kmeans_vs_hc}
  \renewcommand{\arraystretch}{1.2}
  \begin{tabular}{@{}lp{5.5cm}p{5.5cm}@{}}
    \toprule
    \textbf{Feature} & \textbf{K-Means} & \textbf{Hierarchical (Ward)} \\ \midrule
    \textbf{Time Complexity} & $O(N \cdot k \cdot I \cdot d)$ (Linear) & $O(N^2 \log N)$ to $O(N^3)$ \\
    \textbf{Space Complexity} & $O(N \cdot d + k \cdot d)$ (Linear) & $O(N^2)$ (Quadratic) \\
    \textbf{Geometry} & Assumes spherical, convex clusters & Flexible, based on connectivity \\
    \textbf{Parameters} & Requires $k$ beforehand & $k$ determined by cutting tree \\
    \textbf{Robustness} & Sensitive to noise/outliers & More robust (structure-based) \\
    \textbf{Deterministic} & No (Random initialization) & Yes \\ \bottomrule
  \end{tabular}
  \source{Compiled by the author}
\end{table}

\section{Data Generation Framework}
\label{sec:data-generation-framework}

To rigorously evaluate algorithm performance, we implemented a reverse-engineering pipeline that generates a ``Ground Truth'' (Real) dataset and subsequently creates a Synthetic version using the method under investigation.

\subsection{Reference Data Generation (Real)}
The reference datasets were generated using the \texttt{mvtnorm} package in R \cite{mvtnormmanual, genzbretz2009} to simulate controlled multivariate clusters. We varied four key parameters to create a total of 180 unique data configurations. Each configuration was then repeated $M = 100$ times; in the experiment JSON reported in Appendix~\ref{sec:app-json}, this replicate count is stored under the field \texttt{m}:
\begin{itemize}
  \item \textbf{Sample Size ($N$):} Ranging from 100 to 1,000 observations to test scalability.
  \item \textbf{Number of Clusters ($k$):} Set to 2, 3, or 4 to simulate varying levels of complexity.
  \item \textbf{Cluster Separation:} A separation index controlling the overlap between clusters. Higher separation implies distinct, easy-to-detect clusters; lower separation implies significant overlap.
  \item \textbf{Probability Distribution:}
  \begin{itemize}
    \item \textbf{Normal:} Standard Gaussian clusters (spherical).
    \item \textbf{Gamma:} Skewed distributions (non-spherical) to test the robustness of the algorithms against non-normal geometries \cite{yan2007}.
  \end{itemize}
\end{itemize}

\subsection{Synthetic Data Generation (CART Method)}
For each real dataset, a corresponding synthetic dataset was generated using the \textbf{Classification and Regression Trees (CART)} method, implemented in the \texttt{synthpop} R package \cite{nowok2016}.

For the qualitative visualization figures presented later in Chapter~\ref{chap:results}, one additional clarification is necessary. The leftmost synthetic panels reuse simulation-side cluster colouring only as a diagnostic device to compare the generated point cloud with the known reference geometry. These colours do \textbf{not} imply a one-to-one correspondence between original and synthetic records, and such ``true labels'' would not be observable in a realistic synthetic-data release. Accordingly, those panels are interpreted only as simulated diagnostics; the substantive evaluation relies on the clustering predictions and quantitative metrics reported below.

\textbf{The Sequential Regression Approach}

The CART method synthesizes data variable by variable. Let $X = (X_1, X_2, ..., X_p)$ be the vector of variables in the real dataset. The joint distribution is approximated as a product of conditional distributions:
\begin{equation}
  f(X_1, ..., X_p) \approx f(X_1) \cdot f(X_2|X_1) \cdot f(X_3|X_1, X_2) \cdot ... \cdot f(X_p|X_1, ..., X_{p-1})
\end{equation}

The synthesis proceeds sequentially:
\begin{enumerate}
  \item \textbf{Step 1:} Generate $X_1$ by resampling from the original univariate distribution.
  \item \textbf{Step 2:} Train a Decision Tree (CART) to predict $X_2$ using $X_1$ as the predictor.
  \item \textbf{Step 3:} Generate synthetic values for $X_2$ by passing the synthetic $X_1$ through the tree and sampling from the appropriate leaf node.
  \item \textbf{Step 4:} Repeat for all subsequent variables, using all previously generated variables as predictors.
\end{enumerate}

\textbf{Structural Artifacts}

While CART is non-parametric and can capture non-linear relationships, it introduces specific geometric artifacts. Because decision trees split data using orthogonal hyperplanes (rules like $x > 5$), they tend to approximate smooth diagonal curves with ``staircase'' or ``blocky'' boundaries. A key objective of this study is to determine if this ``blocky'' structure causes K-Means or Hierarchical Clustering to fail.

\section{Evaluation Metrics}
\label{sec:evaluation-metrics-methods}

Standard clustering accuracy metrics (like purity) are insufficient when the cluster labels are unknown or permuted. We introduced a set of structural fidelity metrics to quantify the degradation between Real and Synthetic data. Table~\ref{tab:clustering_metrics} provides a summary of the metrics employed.

\begin{table}[H]
  \centering
  \caption[Summary of Clustering Evaluation Metrics]{Summary of the geometric and detection metrics used to evaluate clustering utility.}
  \label{tab:clustering_metrics}
  \renewcommand{\arraystretch}{1.3}
  \begin{tabular}{@{}lp{8cm}l@{}}
    \toprule
    \textbf{Metric} & \textbf{Description} & \textbf{Target} \\ \midrule
    \textbf{Success Rate} & Binary measure of whether the algorithm detected the true $k$. & $1.0$ (100\%) \\
    \textbf{Gini Coefficient} & Measures the inequality of cluster sizes (imbalance). & $\approx 0.0$ \\
    \textbf{MCD} & Mean Centroid Displacement: shift in cluster centers. & $\approx 0.0$ \\
    \textbf{MVD} & Mean Variance Difference: change in cluster dispersion. & $\approx 0.0$ \\ \bottomrule
  \end{tabular}
  \source{Compiled by the author}
\end{table}

\subsection{Preprocessing Step: Cluster Alignment (Hungarian Algorithm)}
It is crucial to note that cluster alignment is not an evaluation metric \textit{per se}, but rather a mandatory preprocessing step required before any comparative structural fidelity metrics can be applied. Because clustering algorithms are fundamentally unsupervised, the cluster labels they assign are arbitrary (e.g., ``Cluster 1'' in the Original Data might be labeled ``Cluster 3'' in the Synthetic Data).

To calculate accuracy or geometric distortion between the two datasets, we must first find the optimal 1-to-1 mapping between the original labels and the synthetic labels. We achieve this by formulating it as the \textbf{Linear Sum Assignment Problem} and solving it using the \textbf{Hungarian Algorithm} \cite{kuhn1955}.

We construct a cost matrix $C$ where $C_{ij}$ represents the number of mismatches if Predicted Cluster $i$ is assigned to True Cluster $j$. The algorithm finds the optimal permutation of labels that minimizes this cost (maximizes the overlap), allowing the subsequent metrics (like MCD and MVD) to be calculated on correctly matched pairs.

\subsection{Success Rate (Detection)}
The primary metric is binary: \textbf{Did the algorithm find the correct number of clusters?} We utilize the \textit{Silhouette Score and the Elbow Method} (for K-Means) \cite{rousseeuw1987} or the \textit{Dendrogram Cut} (for Hierarchical) to estimate $\hat{k}$.
\begin{equation}
  \text{Success} =
  \begin{cases}
    1 & \text{if } \hat{k}_{\text{predicted}} = k_{\text{true}} \\
    0 & \text{otherwise}
  \end{cases}
\end{equation}

\subsection{Structural Fidelity Metrics}
To measure the geometric distortion introduced by synthesis, we calculate the following comparative metrics between the Real and Synthetic clusters:

\textbf{1. Gini Coefficient (Cluster Balance)}

The Gini coefficient measures the inequality of cluster sizes. If the synthetic data generation collapses a small cluster into a large one, the Gini coefficient will increase.
\begin{equation}
  G = \frac{\sum_{i=1}^{k} \sum_{j=1}^{k} |n_i - n_j|}{2k \sum_{i=1}^{k} n_i}
\end{equation}
Where $n_i$ is the number of points in cluster $i$. A value of 0 indicates perfectly equal cluster sizes; a value near 1 indicates high imbalance. To evaluate the structural distortion introduced by the synthesis process, we compute the absolute difference between the Gini coefficient of the Synthetic Data ($G^{(S)}$) and the Original Data ($G^{(R)}$):
\begin{equation}
  \Delta G = |G^{(S)} - G^{(R)}|
\end{equation}
A $\Delta G$ near $0$ indicates that the synthetic data has successfully preserved the original cluster balance.

\textbf{2. Mean Centroid Displacement (MCD)}

We measure how far the centers of the clusters have shifted during synthesis. Let $\mu_j^{(R)}$ be the centroid of cluster $j$ in Real data and $\mu_j^{(S)}$ in Synthetic data:
\begin{equation}
  \text{MCD} = \frac{1}{k} \sum_{j=1}^{k} ||\mu_j^{(R)} - \mu_j^{(S)}||_2
\end{equation}
High displacement indicates that the synthetic data has drifted from the original spatial location.

\textbf{3. Mean Variance Difference (MVD)}

We measure if the synthetic clusters are more dispersed (noisy) or more compact than the real ones. We compare the average within-cluster variance ($\sigma^2$):
\begin{equation}
  \text{MVD} = \frac{1}{k} \sum_{j=1}^{k} \left| \text{Var}(C_j^{(R)}) - \text{Var}(C_j^{(S)}) \right|
\end{equation}
